\SetPageTheme{challenge}
\PageHeader
  {Provocarea redactiei}
  {Trimite propria versiune alterată de TicTacToe, împreună cu explicația felului în care ai schimbat jocul.}
  {Final 1 din 2}

\begin{bossbox}[Ce trimiti]
\begin{enumerate}
\item codul sursa al variantei tale;
\item explicatie (5--10 randuri) despre ce ai schimbat;
\item un output exemplu sau capturi;
\item nume, clasa, titlul versiunii tale.
\end{enumerate}
\end{bossbox}

\begin{solobox}[Checklist inainte de trimitere]
\begin{itemize}
\item Codul ruleaza fara blocaje?
\item Mutarile invalide sunt refuzate corect?
\item Castigul/remiza sunt detectate corect?
\item Modificarile tale sunt explicate clar?
\end{itemize}
\AnswerSpace{6}
\end{solobox}

\begin{conceptbox}[Date de trimitere]
\FieldLine{Email redactiei}
\FieldLine{Formular site}
\FieldLine{Data limita}
\end{conceptbox}

\clearpage

\SetPageTheme{story}
\thispagestyle{empty}
\begin{center}
{\Huge\bfseries Felicitari!}\par
\vspace{1.2em}
{\Large Ai construit drumul complet:}\par
\vspace{0.5em}
{\large stare \(\rightarrow\) decizie \(\rightarrow\) repetitie \(\rightarrow\) grila \(\rightarrow\) joc}
\end{center}

\vspace{1.8em}
\begin{missionbox}[Mesaj final]
Calculatorul nu stie sa numere singur. Dar acum stii tu sa-l inveti:
\begin{itemize}
\item sa tina minte stari;
\item sa actualizeze corect valori;
\item sa aleaga drumuri prin conditii;
\item sa repete procese controlat;
\item sa randeze o tabla de joc si sa aplice reguli de victorie.
\end{itemize}
\end{missionbox}

\vfill
\begin{center}
\QRCodePlaceholder

\smallskip
\small QR rezervat pentru pagina editiei si simulatorul central al acestui numar.
\end{center}

\clearpage
